%!TEX root = ../main.tex

\chapter{Graph RAG Variants}\label{ch:graphrag}

Chapter~\ref{ch:graphhelp} showed two opposite effects of the graph facts. They help when
passage retrieval misses the gold section and a fact still comes from it (4.6 points on
342 validation questions), and they hurt when the passages are already right and the
facts come from other sections (1.0 point on 1{,}005 questions). This chapter tries to
keep the first effect and remove the second, by changing which facts are given to the
generator and by using the graph to rerank the sections. All variants use the
conversational generator (\texttt{gpt-5.4-nano}, G2 + abstain(G0)) and the same rewrites,
passages and linked entities. They were developed on the validation split, with every
setting tuned on dev-500 and reported on the other 3{,}927 questions, and with a stopping
rule fixed in advance. The two subsets above are kept as defined for E and are not
recomputed per variant. Only the systems retained at the end were run, once, on the test
split. As in Chapter~\ref{ch:graphhelp}, P denotes passages and E passages with up to 30
facts.


% ===========================================================================
\section{Variants}\label{sec:graphrag:variants}

A question has on average 47 candidate facts (median 39, 95th percentile 120), the
relations of its ten linked entities; the 30 facts of E take about 1{,}670 prompt tokens.
Four variants and one reference were run:
\begin{description}
    \item[V0] graph facts only (B) with the dialogue history, which was missing so far;
    \item[V1] every candidate fact (relation and evidence quote) scored against the
        rewritten query with the cross-encoder of the passage pipeline; the top $k$ facts
        are kept ($k = 5$ and $k = 10$), and facts below a score threshold are dropped. The
        threshold was searched on dev-500; every threshold lowered \fone{}, so none is used.
        The control is the same number of random facts from the same domain;
    \item[V2] as V1 with $k = 10$, but the slots are filled first with facts whose evidence
        section is not among the five passages: all ten slots (V2a) or five such facts and
        five from the passages (V2b);
    \item[V3] the graph as a reranker signal: each candidate section gets the score
        $\mathrm{ce} + w \cdot g$, with ce the cross-encoder score and $g$ either the best
        cross-encoder score of a fact from that section or the entity-overlap score of A,
        both min-max normalised per question. The candidates are the 50 sections of the
        passage ranking and the evidence sections of the candidate facts (11.6 more per
        question);
    \item[V4] two-hop paths between linked entities, to be run only if V1 to V3 left a
        clear gap.
\end{description}


% ===========================================================================
\section{Results on Validation}\label{sec:graphrag:val}

\begin{table}[ht]
    \centering
    \caption{Graph RAG variants on the full validation split (\%, $n = 4{,}427$). $\Delta$:
    paired bootstrap difference to E with its 95\% interval. Miss / hit: the 342 questions
    where passage retrieval misses the gold section and one of E's facts comes from it, and
    the 1{,}005 where the gold section is in the top 5 and none of E's facts comes from it.
    Tokens: generator input per turn.}
    \label{tab:graphrag:val}
    \small
    \begin{tabular}{lrrlrrr}
        \toprule
        System & \fone{} & answ. & $\Delta$\fone{}, all & Miss & Hit & Tokens \\
        \midrule
        P: passages                    & 27.5 & 27.6 & $-0.5$ ($-0.9$, $-0.1$) & 20.6 & 31.9 & 2{,}481 \\
        E: passages + 30 facts         & 28.0 & 28.2 & ---                     & 25.2 & 30.9 & 5{,}775 \\
        V0: graph facts only           & 24.6 & 24.1 & $-3.4$ ($-3.9$, $-2.8$) & 26.5 & 22.3 & 3{,}741 \\
        V1, top 5 facts                & 27.6 & 28.1 & $-0.4$ ($-0.8$, $-0.0$) & 22.7 & 30.7 & 3{,}162 \\
        V1, 5 random facts             & 27.8 & 27.3 & $-0.2$ ($-0.7$, $0.2$)  & 21.6 & 31.7 & 3{,}089 \\
        V1, top 10 facts               & 28.0 & 28.3 & $-0.0$ ($-0.4$, $0.4$)  & 23.2 & 30.8 & 3{,}778 \\
        V1, 10 random facts            & 28.1 & 27.6 & $+0.1$ ($-0.4$, $0.5$)  & 21.5 & 32.1 & 3{,}713 \\
        V2a, out-of-context first      & 27.8 & 28.0 & $-0.2$ ($-0.6$, $0.2$)  & 25.5 & 30.6 & 3{,}753 \\
        V2b, 5 + 5                     & 27.9 & 28.2 & $-0.1$ ($-0.4$, $0.3$)  & 23.2 & 30.9 & 3{,}770 \\
        \bottomrule
    \end{tabular}
\end{table}

Table~\ref{tab:graphrag:val} shows the generation results. With dialogue history, graph
facts alone reach 24.6\%, half a point more than without it, and stay far below E. Keeping
the ten facts with the best cross-encoder score (V1) gives the same \fone{} as E with 35\%
fewer input tokens, and on answerable questions it is 0.7 points better than ten random
facts (interval $[0.4, 0.9]$); with five facts the model abstains less and loses 0.4 points
on all questions, so $k = 10$ was kept as the default. However, V1 loses 2.0 points against
E on the 342 questions it was meant to keep: the cross-encoder ranks first the facts that
repeat the passages, and the facts from other sections that recovered the retrieval misses
are cut. On the 1{,}005 questions where the passages are right, removing off-topic facts did
not recover the point that E loses there.

Putting the facts from other sections first (V2a) recovers these questions, 2.4 points
over V1 (interval $[1.0, 3.9]$), back to E's level, without a significant loss on all
questions; on answerable questions in general it is 0.3 points below V1, because facts from
the same sections as the passages helped a little when the passages already held the
answer. V2b selects almost the same facts as V1, since V1's top ten already come mostly from
other sections, and changes nothing.

The reranker signal of V3 does not help. On dev-500, recall@5 is 67.8\% with the
cross-encoder alone ($w = 0$) and falls for every positive weight, to 54.6\% at $w = 2$ with
the fact score and 52.2\% with the entity overlap, so the tuned weight is 0 and no generation
run was made. Added to the cross-encoder, the graph score mostly promotes sections that
mention the linked entities without holding the answer. On the 342 questions of the first
subset it does raise recall@5, from 3.5\% to 48.0\% at $w = 2$, but this subset is selected by
a fact from the gold section, which is what the fact score rewards, and it cannot be
identified before generation (Section~\ref{sec:graphhelp:router}).

V4 was not run. Among the 1{,}605 validation questions where passage retrieval misses, 1{,}115
have no one-hop fact that reaches a gold section, and a two-hop path between two linked
entities reaches one in none of them: both edges of such a path are themselves one-hop facts
of a linked entity. A broader two-hop neighbourhood (relations one step beyond the
neighbours of the linked entities) would reach a gold section in 16.9\% of the misses, but it
is a much larger candidate set and was not tried.


% ===========================================================================
\section{Results on the Test Split}\label{sec:graphrag:test}

P, E and V1 with $k = 10$ were frozen and run once on the test split.

\begin{table}[ht]
    \centering
    \caption{Final systems on the test split (\%, $n = 4{,}335$, 241 unanswerable), all with
    the conversational generator. Abstain: abstention recall / false abstentions. Tokens:
    generator input per turn.}
    \label{tab:graphrag:test}
    \small
    \begin{tabular}{lrrrrr}
        \toprule
        System & \fone{} & answerable & unanswerable & Abstain & Tokens \\
        \midrule
        P: passages             & 27.6 & 27.5 & 28.4 & 27.4 / 6.3 & 2{,}324 \\
        E: passages + 30 facts  & 28.2 & 28.4 & 24.1 & 23.7 / 4.4 & 5{,}621 \\
        V1: passages + 10 facts & 27.8 & 28.1 & 22.3 & 21.6 / 4.5 & 3{,}608 \\
        \bottomrule
    \end{tabular}
\end{table}

E gains 0.6 points over P on all questions (interval $[0.1, 1.0]$) and 0.9 on answerable
ones ($[0.6, 1.2]$), as on validation. V1 is 0.4 points below E on all questions
($[-0.7, 0.0]$) and 0.3 on answerable ones ($[-0.6, -0.0]$), so the tie of validation does not
hold on test; it still gains 0.6 over P on answerable questions with 36\% fewer input tokens
than E. Both graph systems abstain less than P and lose on the unanswerable questions, with
wide intervals. Appending the one-hop facts of the linked entities therefore remains the best
use of the graph found in this work, and the filtered selection is the cheaper option at a
small cost in \fone{}.
